\src{181}
\result{COROLLAIRE}{1} \uline{Soient} $\symup{X}$, $\symup{X'}$ \uline{deux variétés non singulières
localement fermées dans un espace projectif, supposons} $\symup{X'}$ \uline{complète. Soient} $\symup{U}$
\uline{un cycle non singulier sur} $\symup{X\times X'}$, $\symup{a}$ \uline{et} $\symup{b}$ \uline{deux points de} $\symup{X'}$ \uline{tels que}
$\symup{U}$ \uline{coupe transversalement les cycles} $\symup{X\times(a)}$ \uline{et} $\symup{X\times(b)}$, \uline{soient} $\symup{Z}$
\uline{et} $\symup{Z'}$ \uline{les cycles non singuliers sur} $\symup{X}$ \uline{tels que}
$\symup{Z\times(a)=(X\times(a)).U}$, $\symup{Z'\times(b)=(X\times(b)).U}$.
\uline{Alors on a}
\[
  \symup{P_X(Z)=P_X(Z').}
\]

En effet, soit $\symup{f_a:X\longrightarrow X\times X'}$ défini par
$\symup{f_a(x)=(x,a)}$, on a alors en vertu de (4.8) la formule
$\symup{P(Z)=f_a^{*}(P(U))}$, de même $\symup{P(Z')=f_b^{*}(P(U))}$. Or, utilisant la formule de
Künneth
\[
  \symup{H^{*}(X\times X',\underline{\Omega}_{X\times X'}^{*})
  =H^{*}(X,\underline{\Omega}_X^{*})\otimes H^{*}(X',\underline{\Omega}_{X'}^{*})}
\]
et le fait que $\symup{H^{0}(X',\Osh_{X'})}$ est réduit aux scalaires, on obtient
facilement $\symup{f_a^{*}=f_b^{*}}$, d'où la conclusion.

Pour tout $\symup{x\in X}$, $\symup{(x)}$ est une sous-variété non singulière de $\symup{X}$ de
codimension $\symup{n}$, et définit donc un élément $\symup{\varepsilon_x}$ de
$\symup{H^{n}(X,\underline{\Omega}_X^{n})}$. Si $\symup{X}$ est une variété projective non singulière, il
résulte du corollaire 1 que $\symup{\varepsilon_x}$ ne dépend pas du point $\symup{x}$ choisi, on
le note $\symup{\varepsilon_X}$ et on l'appelle \uline{classe fondamentale} de
$\symup{H^{n}(X,\underline{\Omega}_X^{n})}$.

\result{REMARQUE}{} Pour avoir une théorie satisfaisante, il faudrait définir
$\symup{P_X(Z)}$ pour des cycles $\symup{Z}$ quelconques, et prouver le théorème 1 pour une
intersection propre de cycles. (Au moment d'écrire cet exposé, cela n'est pas
encore fait en toute généralité). Admettant que cela est fait, le corollaire 1
devient : si $\symup{Z}$ et $\symup{Z'}$ sont deux cycles algébriquement équivalents, alors
$\symup{P_X(Z)=P_X(Z')}$ (énoncé qui ne semble pas résulter de ce qui précède, même si
$\symup{Z}$ et $\symup{Z'}$ sont non singuliers).

\fsection{\normalfont5}{\normalfont\uline{Le théorème de dualité}.}

Dans ce numéro, $\symup{X}$ désigne une variété projective non singulière de dimension
$\symup{n}$.

\result{THÉORÈME}{2} \uline{La classe fondamentale} $\symup{\varepsilon_X}$ \uline{de}
$\symup{H^{n}(X,\underline{\Omega}_X^{n})}$ \uline{est une base de cet espace vectoriel}.

(La démonstration sera donnée plus bas). Moyennant le théorème précédent on peut
donc identifier $\symup{H^{n}(X,\underline{\Omega}_X^{n})}$ au corps $\symup{\underline{k}}$. Considérons maintenant les
accouplements envisagés au n\textsuperscript{o} 2, qui donnent en particulier un
accouplement

